\documentclass[../DoAn.tex]{subfiles}
\begin{document}

\section{Đặc tả chi tiết các bảng dữ liệu hệ thống}
\label{chapter:db_schema}

Để đảm bảo tính súc tích và tuân thủ giới hạn độ dài của báo cáo, Mục 4.2.3 chỉ trình bày sơ đồ thực thể liên kết (E-R) tổng quát và đặc tả ba bảng dữ liệu cốt lõi nhất. Phần phụ lục này cung cấp đặc tả chi tiết cho toàn bộ mười một bảng quan hệ PostgreSQL còn lại, cấu trúc khóa-giá trị của FoundationDB và quy ước lưu trữ đối tượng nhị phân của SeaweedFS được áp dụng trong hệ thống Dyadia thực tế. Các bảng này tuân thủ các quy ước thiết kế và kiểu dữ liệu chuẩn của hệ thống.

\begin{table}[H]
\centering
\small
\renewcommand{\arraystretch}{1.3}
\begin{tabular}{|c|l|l|p{6.2cm}|}
\hline
\textbf{STT} & \textbf{Tên cột} & \textbf{Kiểu dữ liệu} & \textbf{Mô tả} \\ \hline
1 & id & uuid & Định danh tổ chức (khóa chính) \\ \hline
2 & created\_by & uuid & Người tạo tổ chức (khóa ngoại liên kết bảng \texttt{users}) \\ \hline
3 & name & text & Tên hiển thị của tổ chức \\ \hline
4 & slug & text & Chuỗi định danh duy nhất dùng cho đường dẫn URL \\ \hline
5 & status & organization\_status & Trạng thái hoạt động của tổ chức \\ \hline
6 & created\_at & timestamp & Thời điểm tạo tổ chức \\ \hline
7 & updated\_at & timestamp & Thời điểm cập nhật thông tin tổ chức \\ \hline
\end{tabular}
\caption{Đặc tả các cột của bảng \texttt{organizations}}
\label{tbl:organizations}
\end{table}

\begin{table}[H]
\centering
\small
\renewcommand{\arraystretch}{1.3}
\begin{tabular}{|c|l|l|p{6.2cm}|}
\hline
\textbf{STT} & \textbf{Tên cột} & \textbf{Kiểu dữ liệu} & \textbf{Mô tả} \\ \hline
1 & organization\_id & uuid & Định danh tổ chức (khóa chính, khóa ngoại liên kết bảng \texttt{organizations}) \\ \hline
2 & user\_id & uuid & Định danh người dùng (khóa chính, khóa ngoại liên kết bảng \texttt{users}) \\ \hline
3 & role & organization\_role & Vai trò của người dùng trong tổ chức (quản trị, giảng viên, học viên) \\ \hline
4 & status & member\_status & Trạng thái thành viên trong tổ chức (đã mời, đang hoạt động) \\ \hline
5 & created\_at & timestamp & Thời điểm người dùng gia nhập tổ chức \\ \hline
6 & updated\_at & timestamp & Thời điểm cập nhật trạng thái thành viên \\ \hline
\end{tabular}
\caption{Đặc tả các cột của bảng \texttt{organization\_members}}
\label{tbl:organization_members}
\end{table}

\begin{table}[H]
\centering
\small
\renewcommand{\arraystretch}{1.3}
\begin{tabular}{|c|l|l|p{6.2cm}|}
\hline
\textbf{STT} & \textbf{Tên cột} & \textbf{Kiểu dữ liệu} & \textbf{Mô tả} \\ \hline
1 & id & uuid & Định danh khóa học (khóa chính) \\ \hline
2 & organization\_id & uuid & Tổ chức sở hữu khóa học (khóa ngoại liên kết bảng \texttt{organizations}) \\ \hline
3 & owner\_id & uuid & Giảng viên chủ sở hữu khóa học (khóa ngoại liên kết bảng \texttt{users}) \\ \hline
4 & title & text & Tên khóa học \\ \hline
5 & description & text & Mô tả chi tiết nội dung khóa học \\ \hline
6 & status & course\_status & Trạng thái khóa học (bản nháp, đã xuất bản, đã khóa) \\ \hline
7 & created\_at & timestamp & Thời điểm tạo khóa học \\ \hline
8 & updated\_at & timestamp & Thời điểm cập nhật thông tin khóa học \\ \hline
\end{tabular}
\caption{Đặc tả các cột của bảng \texttt{courses}}
\label{tbl:courses}
\end{table}

\begin{table}[H]
\centering
\small
\renewcommand{\arraystretch}{1.3}
\begin{tabular}{|c|l|l|p{6.2cm}|}
\hline
\textbf{STT} & \textbf{Tên cột} & \textbf{Kiểu dữ liệu} & \textbf{Mô tả} \\ \hline
1 & course\_id & uuid & Định danh khóa học (khóa chính, khóa ngoại liên kết bảng \texttt{courses}) \\ \hline
2 & user\_id & uuid & Định danh người dùng (khóa chính, khóa ngoại liên kết bảng \texttt{users}) \\ \hline
3 & role & course\_role & Vai trò trong khóa học (giảng viên trợ giảng, học viên) \\ \hline
4 & created\_at & timestamp & Thời điểm tham gia khóa học \\ \hline
5 & updated\_at & timestamp & Thời điểm cập nhật vai trò \\ \hline
\end{tabular}
\caption{Đặc tả các cột của bảng \texttt{course\_members}}
\label{tbl:course_members}
\end{table}

\begin{table}[H]
\centering
\small
\renewcommand{\arraystretch}{1.3}
\begin{tabular}{|c|l|l|p{6.2cm}|}
\hline
\textbf{STT} & \textbf{Tên cột} & \textbf{Kiểu dữ liệu} & \textbf{Mô tả} \\ \hline
1 & course\_id & uuid & Định danh khóa học (khóa chính, khóa ngoại liên kết bảng \texttt{courses}) \\ \hline
2 & user\_id & uuid & Học sinh gửi yêu cầu tham gia (khóa chính, khóa ngoại liên kết bảng \texttt{users}) \\ \hline
3 & status & join\_request\_status & Trạng thái yêu cầu (đang chờ duyệt, đã đồng ý, từ chối) \\ \hline
4 & created\_at & timestamp & Thời điểm gửi yêu cầu \\ \hline
5 & updated\_at & timestamp & Thời điểm phê duyệt hoặc từ chối yêu cầu \\ \hline
6 & requested\_role & course\_role & Vai trò đăng ký tham gia khóa học (giảng viên, học viên) \\ \hline
\end{tabular}
\caption{Đặc tả các cột của bảng \texttt{course\_join\_requests}}
\label{tbl:course_join_requests}
\end{table}

\begin{table}[H]
\centering
\small
\renewcommand{\arraystretch}{1.3}
\begin{tabular}{|c|l|l|p{6.2cm}|}
\hline
\textbf{STT} & \textbf{Tên cột} & \textbf{Kiểu dữ liệu} & \textbf{Mô tả} \\ \hline
1 & id & uuid & Định danh học phần (khóa chính) \\ \hline
2 & course\_id & uuid & Khóa học chứa học phần (khóa ngoại liên kết bảng \texttt{courses}) \\ \hline
3 & title & text & Tiêu đề chương học \\ \hline
4 & order\_index & integer & Thứ tự sắp xếp của chương học trong khóa học \\ \hline
5 & created\_at & timestamp & Thời điểm tạo học phần \\ \hline
6 & updated\_at & timestamp & Thời điểm cập nhật tên học phần \\ \hline
\end{tabular}
\caption{Đặc tả các cột của bảng \texttt{course\_modules}}
\label{tbl:course_modules}
\end{table}

\begin{table}[H]
\centering
\small
\renewcommand{\arraystretch}{1.3}
\begin{tabular}{|c|l|l|p{6.2cm}|}
\hline
\textbf{STT} & \textbf{Tên cột} & \textbf{Kiểu dữ liệu} & \textbf{Mô tả} \\ \hline
1 & id & uuid & Định danh phân đoạn (khóa chính) \\ \hline
2 & lesson\_id & uuid & Bài giảng tương ứng (khóa ngoại liên kết bảng \texttt{lessons}) \\ \hline
3 & order\_index & integer & Thứ tự của phân đoạn trong bài học \\ \hline
4 & start\_seconds & float & Thời điểm bắt đầu của phân đoạn trong video (giây) \\ \hline
5 & end\_seconds & float & Thời điểm kết thúc của phân đoạn trong video (giây) \\ \hline
6 & summary & text & Văn bản tóm tắt ý chính của phân đoạn \\ \hline
7 & question\_count\_config & integer & Cấu hình số lượng câu hỏi cần sinh tự động cho phân đoạn này \\ \hline
8 & interaction\_config & jsonb & Cấu hình chi tiết sinh câu hỏi riêng cho phân đoạn \\ \hline
9 & created\_at & timestamp & Thời điểm tạo phân đoạn \\ \hline
10 & updated\_at & timestamp & Thời điểm chỉnh sửa phân đoạn \\ \hline
\end{tabular}
\caption{Đặc tả các cột của bảng \texttt{lesson\_transcript\_chunks}}
\label{tbl:lesson_transcript_chunks}
\end{table}

\begin{table}[H]
\centering
\small
\renewcommand{\arraystretch}{1.3}
\begin{tabular}{|c|l|l|p{6.2cm}|}
\hline
\textbf{STT} & \textbf{Tên cột} & \textbf{Kiểu dữ liệu} & \textbf{Mô tả} \\ \hline
1 & id & uuid & Định danh lượt làm bài (khóa chính) \\ \hline
2 & user\_id & uuid & Học viên thực hiện làm bài (khóa ngoại liên kết bảng \texttt{users}) \\ \hline
3 & lesson\_id & uuid & Bài giảng tương ứng (khóa ngoại liên kết bảng \texttt{lessons}) \\ \hline
4 & total\_score & real & Điểm số thực tế học viên đạt được \\ \hline
5 & max\_score & real & Điểm số tối đa của toàn bộ bài học tương tác \\ \hline
6 & status & text & Trạng thái làm bài (đang làm bài, đã nộp bài) \\ \hline
7 & submitted\_at & timestamptz & Thời điểm nộp bài \\ \hline
8 & attempt\_count & integer & Thứ tự lượt làm bài của học viên đối với bài học \\ \hline
9 & started\_at & timestamptz & Thời điểm bắt đầu làm bài giảng \\ \hline
10 & video\_watch\_fraction & real & Tỷ lệ thời lượng video mà học viên thực sự đã xem \\ \hline
\end{tabular}
\caption{Đặc tả các cột của bảng \texttt{lesson\_attempts}}
\label{tbl:lesson_attempts}
\end{table}

\begin{table}[H]
\centering
\small
\renewcommand{\arraystretch}{1.3}
\begin{tabular}{|c|l|l|p{6.2cm}|}
\hline
\textbf{STT} & \textbf{Tên cột} & \textbf{Kiểu dữ liệu} & \textbf{Mô tả} \\ \hline
1 & attempt\_id & uuid & Lượt làm bài tương ứng (khóa chính, khóa ngoại liên kết bảng \texttt{lesson\_attempts}) \\ \hline
2 & interaction\_id & uuid & Định danh câu hỏi tương tác (khóa chính, khóa ngoại liên kết bảng \texttt{lesson\_interactions}) \\ \hline
3 & response & jsonb & Dữ liệu câu trả lời của học viên (lưu động theo loại câu hỏi) \\ \hline
4 & score & real & Điểm số đạt được cho câu trả lời này \\ \hline
5 & max\_score & real & Điểm tối đa của câu hỏi này \\ \hline
6 & feedback & text & Nhận xét phản hồi tự động từ hệ thống hoặc AI cho câu trả lời \\ \hline
7 & time\_to\_answer\_ms & integer & Thời gian phản hồi/trả lời câu hỏi (mili-giây) \\ \hline
8 & replay\_count & integer & Số lần học viên tua/xem lại đoạn video chứa câu hỏi \\ \hline
9 & metrics & jsonb & Các chỉ số đo lường chi tiết hành vi làm bài của học viên \\ \hline
\end{tabular}
\caption{Đặc tả các cột của bảng \texttt{lesson\_attempt\_responses}}
\label{tbl:lesson_attempt_responses}
\end{table}

\begin{table}[H]
\centering
\small
\renewcommand{\arraystretch}{1.3}
\begin{tabular}{|c|l|l|p{6.2cm}|}
\hline
\textbf{STT} & \textbf{Tên cột} & \textbf{Kiểu dữ liệu} & \textbf{Mô tả} \\ \hline
1 & id & uuid & Định danh tác vụ (khóa chính) \\ \hline
2 & lesson\_id & uuid & Bài giảng gắn liền với tác vụ (khóa ngoại liên kết bảng \texttt{lessons}) \\ \hline
3 & chunk\_id & uuid & Phân đoạn tương ứng được xử lý (nếu có) \\ \hline
4 & task\_type & text & Loại tác vụ (ví dụ: trích xuất audio, phiên âm, sinh câu hỏi) \\ \hline
5 & status & task\_status & Trạng thái tác vụ (đang chờ, đang xử lý, thành công, thất bại, đã hủy) \\ \hline
6 & worker\_id & uuid & Định danh của tiến trình nền nhận thực hiện tác vụ \\ \hline
7 & heartbeat & timestamptz & Thời gian gửi tín hiệu hoạt động cuối cùng của worker \\ \hline
8 & error\_msg & text & Nội dung thông báo lỗi chi tiết nếu tác vụ thất bại \\ \hline
9 & input\_payload & bytea & Dữ liệu nhị phân đầu vào cấu hình cho tác vụ \\ \hline
10 & output\_payload & bytea & Kết quả nhị phân đầu ra sau khi xử lý thành công \\ \hline
11 & queue\_seq & bigint & Số thứ tự hàng đợi để ưu tiên xử lý trước sau \\ \hline
12 & created\_at & timestamptz & Thời điểm tạo tác vụ \\ \hline
13 & updated\_at & timestamptz & Thời điểm cập nhật trạng thái tác vụ gần nhất \\ \hline
14 & started\_at & timestamptz & Thời điểm bắt đầu thực thi tác vụ \\ \hline
15 & finished\_at & timestamptz & Thời điểm kết thúc hoặc thất bại tác vụ \\ \hline
16 & created\_by & uuid & Người khởi tạo tác vụ (khóa ngoại liên kết bảng \texttt{users}) \\ \hline
17 & progress\_step & text & Tên bước tiến trình đang thực thi hiện tại \\ \hline
18 & progress\_current & integer & Chỉ số tiến độ hiện tại \\ \hline
19 & progress\_total & integer & Tổng số bước tiến độ dự kiến \\ \hline
20 & message & text & Thông điệp tiến trình bổ sung hiển thị cho người dùng \\ \hline
\end{tabular}
\caption{Đặc tả các cột của bảng \texttt{tasks}}
\label{tbl:tasks}
\end{table}

\begin{table}[H]
\centering
\small
\renewcommand{\arraystretch}{1.3}
\begin{tabular}{|c|l|l|p{6.2cm}|}
\hline
\textbf{STT} & \textbf{Tên cột} & \textbf{Kiểu dữ liệu} & \textbf{Mô tả} \\ \hline
1 & id & uuid & Định danh trạng thái phân tích (khóa chính) \\ \hline
2 & lesson\_id & uuid & Bài giảng tương ứng (khóa ngoại liên kết bảng \texttt{lessons}) \\ \hline
3 & status & lesson\_analysis\_status & Trạng thái phân tích (pending, processing, done, error, transcript\_extracted, chunks\_ready) \\ \hline
4 & error\_msg & text & Chi tiết thông báo lỗi nếu phân tích thất bại \\ \hline
5 & created\_at & timestamp & Thời điểm khởi tạo phân tích \\ \hline
6 & updated\_at & timestamp & Thời điểm cập nhật trạng thái gần nhất \\ \hline
\end{tabular}
\caption{Đặc tả các cột của bảng \texttt{lesson\_analyses}}
\label{tbl:lesson_analyses}
\end{table}

\subsection{Thiết kế cơ sở dữ liệu khóa-giá trị FoundationDB}
Hệ thống sử dụng FoundationDB (kho lưu trữ khóa-giá trị NoSQL) để lưu trữ bản chép lời video, thông tin phân đoạn nội dung và lược đồ theo dõi tiến trình xem video (watch coverage) theo thời gian thực. Cấu trúc các khóa và không gian tên (namespaces) được đặc tả chi tiết ở Bảng~\ref{tbl:fdb_schema}.

\begin{table}[H]
\centering
\small
\renewcommand{\arraystretch}{1.3}
\begin{tabular}{|c|l|>{\raggedright\arraybackslash}p{4.2cm}|>{\raggedright\arraybackslash}p{5.5cm}|}
\hline
\textbf{STT} & \textbf{Không gian tên} & \textbf{Cấu trúc khóa (Key)} & \textbf{Mô tả dữ liệu giá trị (Value)} \\ \hline
1 & lesson & \texttt{("lesson",\allowbreak "m",\allowbreak lesson\_id,\allowbreak "transcript")} & Số nguyên 4-byte lưu số lượng khối dữ liệu. \\ \hline
2 & lesson & \texttt{("lesson",\allowbreak "c",\allowbreak lesson\_id,\allowbreak "transcript",\allowbreak idx)} & Khối nhị phân chứa nội dung bản chép lời (\texttt{FdbTranscript} Protobuf). \\ \hline
3 & lesson & \texttt{("lesson",\allowbreak "m",\allowbreak lesson\_id,\allowbreak "segments")} & Số nguyên 4-byte lưu số lượng khối phân đoạn. \\ \hline
4 & lesson & \texttt{("lesson",\allowbreak "c",\allowbreak lesson\_id,\allowbreak "segments",\allowbreak idx)} & Khối nhị phân chứa danh sách phân đoạn (\texttt{FdbSegmentsBlob} Protobuf). \\ \hline
5 & chunk & \texttt{("chunk",\allowbreak "m",\allowbreak chunk\_id,\allowbreak "transcript")} & Số nguyên 4-byte lưu số lượng khối phân đoạn. \\ \hline
6 & chunk & \texttt{("chunk",\allowbreak "c",\allowbreak chunk\_id,\allowbreak "transcript",\allowbreak idx)} & Khối nhị phân chứa bản chép lời phân đoạn (\texttt{FdbTranscript} Protobuf). \\ \hline
7 & watch\_cov & \texttt{("watch\_cov",\allowbreak "r",\allowbreak user\_id,\allowbreak lesson\_id)} & Mảng byte đại diện cho bitmap giây xem video. \\ \hline
8 & watch\_cov\_ts & \texttt{("watch\_cov\_ts",\allowbreak "r",\allowbreak user\_id,\allowbreak lesson\_id)} & Số nguyên 8-byte lưu Unix timestamp của lượt xem đầu tiên. \\ \hline
\end{tabular}
\caption{Đặc tả cấu trúc khóa và giá trị trong FoundationDB}
\label{tbl:fdb_schema}
\end{table}

\subsection{Thiết kế kho lưu trữ đối tượng SeaweedFS}
Hệ thống sử dụng SeaweedFS (S3 compatible object storage) để lưu trữ các tệp tin nhị phân dung lượng lớn bao gồm video bài giảng và các tệp âm thanh ghi âm/TTS. Quy ước các đường dẫn đối tượng được đặc tả chi tiết ở Bảng~\ref{tbl:s3_schema}.

\begin{table}[H]
\centering
\small
\renewcommand{\arraystretch}{1.3}
\begin{tabular}{|c|>{\raggedright\arraybackslash}p{4.8cm}|l|>{\raggedright\arraybackslash}p{5.2cm}|}
\hline
\textbf{STT} & \textbf{Đường dẫn đối tượng (Key)} & \textbf{Định dạng} & \textbf{Mô tả} \\ \hline
1 & \texttt{lessons/\allowbreak <lesson\_id>/\allowbreak video.mp4} & mp4 & Tệp video bài giảng gốc do giảng viên tải lên. \\ \hline
2 & \texttt{lessons/\allowbreak <lesson\_id>/\allowbreak ai-audio/\allowbreak <uuid>.wav} & wav & Tệp âm thanh tương tác sinh tự động bằng công nghệ TTS (Piper). \\ \hline
3 & \texttt{lessons/\allowbreak <lesson\_id>/\allowbreak ai-audio-preview/\allowbreak <uuid>.wav} & wav & Tệp âm thanh nghe thử bài nghe do giảng viên sinh ra trong khâu biên tập. \\ \hline
4 & \texttt{lessons/\allowbreak <lesson\_id>/\allowbreak student-recordings/\allowbreak <filename>} & wav, webm & Tệp ghi âm câu trả lời phát âm thực tế của học viên để chấm điểm. \\ \hline
5 & \texttt{seed/\allowbreak <org-slug>/\allowbreak <path>} & mp4, wav, png & Tệp dữ liệu mẫu ban đầu của hệ thống. \\ \hline
\end{tabular}
\caption{Quy ước đường dẫn đối tượng trong Bucket SeaweedFS}
\label{tbl:s3_schema}
\end{table}

\end{document}
